\documentclass[11pt]{amsart}
\usepackage{amsmath,amssymb,amsthm,mathtools}
\usepackage[margin=1.1in]{geometry}
\usepackage{booktabs,array}
\usepackage[hidelinks]{hyperref}
\newtheorem{theorem}{Theorem}[section]
\newtheorem{lemma}[theorem]{Lemma}
\newtheorem{proposition}[theorem]{Proposition}
\theoremstyle{remark}
\newtheorem{remark}[theorem]{Remark}
\numberwithin{equation}{section}
\newcommand{\R}{\mathbb R}
\newcommand{\C}{\mathbb C}
\newcommand{\Z}{\mathbb Z}
\newcommand{\D}{\mathbb D}
\newcommand{\T}{\mathbb T}
\newcommand{\F}{\mathcal F}
\newcommand{\Ad}{\operatorname{Ad}}
\newcommand{\re}{\operatorname{Re}}
\newcommand{\im}{\operatorname{Im}}
\newcommand{\norm}[1]{\lVert#1\rVert}
\newcommand{\abs}[1]{\lvert#1\rvert}
\newcommand{\kf}{\mathfrak k}
\newcommand{\tf}{\mathfrak t}
\setlength{\emergencystretch}{1em}
\newcommand{\Pf}{P_f}
\newcommand{\Pt}{P_t}
\newcommand{\rFourHash}{e3479208ba7234e291cdde4d30e8e8c66ece07ae3bdade67daf9661f622a6337}
\newcommand{\rFourPayload}{77cebdd65d09b3f10ab95a17efb05af1f5b8add1837c1a5d6973b679ced1ba97}
\newcommand{\rFourSeconds}{9.514}
\newcommand{\rFourFullSeconds}{374.957}
\newcommand{\rFourteenHash}{3ea6375202d31df994cb407603dc7d3c7cbee83587b95c91ab4fe71e8a3b5702}
\newcommand{\rFourteenPayload}{a6229049ed6e91c4a690c1f0ed9ad49034272328751d70207cc5bda6829dc88e}
\newcommand{\rFourteenSeconds}{34.589}
\newcommand{\rFourteenFullSeconds}{518.356}
\hypersetup{pdftitle={Counterexamples to Berenstein's conjecture in dimensions four and fourteen},pdfauthor={Jizhou Guo}}
\begin{document}
\title{Counterexamples to Berenstein's conjecture in dimensions four and fourteen}
\author{Jizhou Guo}
\email{mitsuha2021b@gmail.com}
\date{October 10, 2026}
\subjclass[2020]{Primary 35N25; Secondary 35P05, 35J05, 22E30, 65G20}
\keywords{Berenstein conjecture, overdetermined Dirichlet problem, sign-changing eigenfunction, convex domain, interval arithmetic}
\begin{abstract}
There exists a bounded strictly convex domain in $\R^4$, different from a
ball, with real-analytic boundary diffeomorphic to $S^3$, carrying a
real-analytic sign-changing solution of $\Delta u+u=0$ with $u=0$ and
$\partial_\nu u=1$ on the boundary. There also exists a bounded convex
non-ball domain in $\R^{14}$, invariant under the adjoint action of $G_2$
and with strictly convex Cartan section, carrying a sign-changing solution
with zero Dirichlet data and nonzero constant normal derivative.
The constructions combine planar reductions, a conformal formulation on
the unit disc and interval bounds for an infinite-dimensional contraction.
\end{abstract}
\maketitle
\raggedbottom

\section{Introduction}
Let $\nu$ denote the outward unit normal. We consider the overdetermined
Dirichlet eigenvalue problem
\begin{equation}\label{eq:problem}
 \Delta u+\lambda u=0\quad\hbox{in }\Omega,\qquad
 u=0,\quad \partial_\nu u=c\quad\hbox{on }\partial\Omega,
 \qquad \lambda>0,\quad c\ne0.
\end{equation}
Berenstein's conjecture asks whether a bounded smooth domain admitting
such a solution must be a ball. The inverse spectral background is described by Berenstein~\cite{Berenstein1980}; related questions for the Neumann problem were studied by Berenstein--Yang~\cite{BerensteinYang1987}. A dilation normalizes
$\lambda$ to one, and multiplication of $u$ by a constant normalizes $c$
to one. For a sign-definite eigenfunction, Serrin's symmetry theorem
forces the domain to be a ball~\cite{Serrin1971}. The functions constructed
below change sign.

The complementary Schiffer problem prescribes constant Dirichlet data
and zero normal derivative for a nonconstant Helmholtz solution.
The Pompeiu problem asks whether a continuous function whose integral
over every rigid motion of a domain vanishes must be zero. Its
Fourier-spectral formulation was developed by Brown, Schreiber and
Taylor~\cite{BrownSchreiberTaylor1973}. For \eqref{eq:problem}, Green's
identity with $e^{i\omega\cdot x}$, $\abs{\omega}^2=\lambda$, gives
$\int_{\partial\Omega}e^{i\omega\cdot x}\,d\sigma=0$. The boundary measure
appears here, whereas the area measure appears in the corresponding
Schiffer formulation.

Colbrook and Stepaniants~\cite{ColbrookStepaniants2026} developed a
conformal fixed-disc construction for the planar Schiffer and Pompeiu
problems. Colbrook, Sadeghi and Stepaniants~\cite{CSS2026} adapted the
method to \eqref{eq:problem}, obtaining a bounded simply connected
planar non-disc with analytic boundary and a sign-changing eigenfunction.
The author's paper~\cite{Guo2026}, version~2, contains constructions for
the Schiffer and Pompeiu problems and a separate part giving a strictly
convex planar domain for Berenstein's problem. Question~2 of \cite{Guo2026} asks whether Berenstein's conjecture fails in every dimension $n\ge3$, and within the convex class in those dimensions.
The two constructions in the present paper address that question in
dimensions four and fourteen.
Dai, Sun, Wei and Zhang~\cite{DaiSunWeiZhang2026} show that a bounded uniformly convex planar domain with connected $C^{2,\epsilon}$ boundary, $\epsilon\in(0,1)$, is a disc if \eqref{eq:problem} admits a nontrivial solution corresponding to a large eigenvalue.

There are also bifurcation results for related sign-changing problems.
Ruiz~\cite{Ruiz2025} constructs a bounded non-ball domain for a suitable
$C^2$ semilinear nonlinearity. Dai and Zhang~\cite{DaiZhang2025},
Theorem~1.3, obtain a branch of bounded-domain solutions of
$-\rho\Delta u=u-(u^+)^3$ from a bifurcation interval, for symmetry groups
satisfying their spectral condition. Minlend~\cite{Minlend2023},
Theorem~1.1, constructs linear Dirichlet examples on unbounded periodic
domains bifurcating from generalized cylinders. These results have
different nonlinearities or different domain classes from the bounded
linear problem considered here.

\begin{theorem}\label{thm:r4}
There exist a bounded strictly convex domain $\Omega\subset\R^4$,
different from a ball, with real-analytic boundary diffeomorphic to $S^3$,
and a real-analytic sign-changing function $u$, analytic on a
neighbourhood of $\overline\Omega$, such that
\[
 \Delta u+u=0\quad\hbox{in }\Omega,\qquad
 u=0,\quad\partial_\nu u=1\quad\hbox{on }\partial\Omega.
\]
\end{theorem}

\begin{theorem}\label{thm:r14}
There exists a bounded convex domain $\Omega\subset\R^{14}$, different
from a ball, carrying a sign-changing solution of
$\Delta u+u=0$ with $u=0$ and nonzero constant normal derivative on its
boundary. After an isometric identification $\R^{14}\simeq\mathfrak g_2$,
the domain is invariant under the adjoint action of the compact group
$G_2$, and its Cartan section is strictly convex.
\end{theorem}

We give the planar reductions and coefficient estimates before specifying
the two contraction balls. Every numerical upper bound used below is
either a rational consequence of those estimates or an interval
comparison in the accompanying verification files. The analytic
arguments explain why the finite calculations control every infinite
column and why an exact coefficient zero gives the stated boundary data.

\section{The planar reductions}\label{sec:reductions}
\subsection{Dimension four}
Write $X=(x,X')\in\R\times\R^3$ and $y=\abs{X'}$. For an
$O(3)$-invariant function $u(x,y)$, direct differentiation gives
\begin{equation}\label{eq:r4radial}
 \Delta_{x,y}(yu)=y\left(u_{xx}+u_{yy}+\frac2y u_y\right)
                 =y\Delta_{\R^4}u,\qquad y>0.
\end{equation}
Thus we seek a plane domain $D$, symmetric under both coordinate
reflections, and an analytic function $F$ satisfying
\begin{equation}\label{eq:r4plane}
 (\Delta+1)F=0,\quad F|_{\partial D}=0,\quad
 \partial_\nu F=y\quad\hbox{on }\partial D,\qquad
 F(x,-y)=-F(x,y).
\end{equation}
Odd analytic division gives $F(x,y)=yG(x,y^2)$ near the axis. Consequently
$u(x,X')=G(x,\abs{X'}^2)$ is analytic across it. On the regular boundary,
the product rule and $u=0$ give $\partial_\nu F=y\partial_\nu u$.
Equation \eqref{eq:r4radial} and continuity then give the equation and
constant normal derivative also at the axis.

\subsection{Dimension fourteen}
Equip the compact Lie algebra $\kf=\mathfrak g_2$ with an
$\Ad(G_2)$-invariant Euclidean inner product, and let $\tf$ be its
two-dimensional Cartan subalgebra. Choose orthonormal coordinates
$z=x+iy$ so that the reflection walls are $\arg z=k\pi/6$. The product of
the six positive roots is a nonzero scalar multiple of
\[
 \pi(z)=\im z^6.
\]
We use this normalization of $\pi$. It is harmonic. Equivalently,
$\Delta\pi$ is an anti-invariant polynomial of degree four, while every
anti-invariant polynomial is divisible by the six wall factors.

On the regular Cartan plane, each root plane has real dimension two.
The orbit-volume density is therefore a constant times $\pi^2$.
For an invariant function $u$, the radial part of the Euclidean
Laplacian is
\begin{equation}\label{eq:radialg2}
 (\Delta_\kf u)|_\tf
 =\pi^{-2}\operatorname{div}_\tf(\pi^2\nabla u)
 =\pi^{-1}\Delta_\tf(\pi u).
\end{equation}
The second equality follows by expanding the derivatives and using
$\Delta\pi=0$. This also follows directly from the orthogonal root-plane
coordinates and their Jacobian. An invariant boundary has its normal in
the Cartan plane at a regular point: its tangent space contains the
orbit tangent space, whose orthogonal complement is $\tf$. Hence, if
$F=\pi u$ and $u=0$, then
\begin{equation}\label{eq:fluxg2}
 \partial_\nu F=\pi\,\partial_\nu u.
\end{equation}
It suffices to construct a strictly convex Weyl-invariant section $D$
and an anti-invariant planar field with
$(\Delta+1)F=0$, $F=0$ and $\partial_\nu F=\pi$ on $\partial D$.
The lift and its singular points are treated in
Section~\ref{sec:reconstruction}.

\section{The fixed disc and coefficient spaces}\label{sec:disc}
We use $m=1$ for the meridian problem in dimension four and $m=6$
for the Cartan problem in dimension fourteen. This use of $m=1$ is a
notation for the Fourier sector, rather than a Lie-algebra root count.
Let
\[
 \psi(w)=\sum_{j\equiv1\ (2m)}c_jw^j,\qquad c_j\in\R,\quad c_1>0,
 \qquad D=\psi(\D).
\]
Take a fixed positive field scale $B$, with $B=1$ in dimension four, and
put $F=B V\circ\psi^{-1}$ and $p=\psi'$. The pulled-back equation is
$\Delta V+\abs p^2V=0$.

\subsection{The full zero-Dirichlet inverse}
For $n\equiv m\pmod{2m}$, $n\ge m$, and $s\ge0$, set
\[
 S_{n,s}(re^{i\theta})
 =r^nP_s^{(0,n)}(2r^2-1)\sin(n\theta).
\]
The Jacobi polynomials are normalized by $P_s^{(0,n)}(1)=1$. The
zero-Dirichlet inverse of $\Delta$ on this sector is
\begin{align}
 K_DS_{n,0}&=\frac{S_{n,1}-S_{n,0}}{4(n+1)(n+2)},\label{eq:kd0}\\
 K_DS_{n,s}&=\frac{S_{n,s-1}}{4d(d+1)}
             -\frac{S_{n,s}}{2d(d+2)}
             +\frac{S_{n,s+1}}{4(d+1)(d+2)},\quad
 d=n+2s,\quad s\ge1.\label{eq:kds}
\end{align}
Substitution in the Jacobi differential equation verifies
$\Delta K_DS_{n,s}=S_{n,s}$. The coefficients of each column sum to zero,
which gives its zero boundary value. Differentiating at $r=1$ gives
\begin{equation}\label{eq:normal}
 Ng:=\partial_rK_Dg|_\T
   =\sum_n\frac{g_{n,0}}{2(n+1)}\sin(n\theta).
\end{equation}
In particular the harmonic source row $s=0$ is retained. Removing it
would remove the required nonzero normal trace. These full Dirichlet
formulas are also used in~\cite{CSS2026}.

Define
\[
 \chi_n(\theta)=\frac{\sin(n\theta)}{\sin(m\theta)}
 =U_{n/m-1}(\cos(m\theta)),\qquad
 q=\frac{Ng}{\sin(m\theta)},\qquad
 h=\frac{\im(\psi^m)}{B\sin(m\theta)}.
\]
The quotients extend analytically to the walls. With $V=K_Dg$, the exact
coefficient equations are
\begin{equation}\label{eq:map}
 \F_1(g,c)=g+\abs p^2K_Dg,\qquad
 \F_2(g,c)=\frac{q^2-h^2\abs p^2}{\eta},
\end{equation}
where $\eta=1$ in dimension four and $\eta=b$ in dimension fourteen;
$b$ is the fixed coefficient $c_1$ at the inverse centre. Thus $\eta$
does not vary with the unknown. The boundary polynomial has degree
$2m+2$. If $q,h>0$ and $p\ne0$ on the boundary, a zero of
\eqref{eq:map} gives $q=h\abs p$, which is the signed normal condition.

\subsection{Norms and multiplication}
For $\rho>1$ define
\begin{align}
 \norm g_G&=\sum_{n,s}\rho^n\abs{g_{n,s}},&
 \norm v_U&=\sum_n\rho^{n-m}\abs{v_n},&
 \norm v_{U,1}&=\sum_n\frac nm\rho^{n-m}\abs{v_n},\label{eq:norms}\\
 \norm c_C&=\sigma\sum_j(j+2m)\rho^{j-1}\abs{c_j},&
 X&=G\oplus C,&
 \mathcal Y&=G\oplus U.\label{eq:shapenorm}
\end{align}
Both direct sums have the sum norm. We take
$(\rho,\sigma)=(11/10,b^3)$ in dimension four and
$(51/50,1)$ in dimension fourteen.

Here are the multiplication estimates used throughout. In the complex
mother space let
$\Phi_{k,s}=r^{|k|}P_s^{(0,|k|)}(2r^2-1)e^{ik\theta}$, with weight
$\rho^{|k|}$. For $k\ge0$,
\[
 w\Phi_{k,s}
 =\frac{k+s+1}{k+2s+1}\Phi_{k+1,s}
  +\frac{s}{k+2s+1}\Phi_{k+1,s-1}.
\]
For $k=-n<0$ the corresponding coefficients are
$(n+s)/(n+2s+1)$ at $(k+1,s)$ and
$(s+1)/(n+2s+1)$ at $(k+1,s+1)$. The conjugate recurrence follows by
conjugation. All coefficients are nonnegative and sum to one. Iteration
therefore gives
\begin{equation}\label{eq:wiener}
 \norm{af}_G\le \norm a_{\mathrm W}\norm f_G,\qquad
 \norm a_{\mathrm W}=\sum_k\rho^{|k|}\abs{a_k}.
\end{equation}
The sine embedding uses conjugate coefficients of magnitude
$|g_{n,s}|/2$, so its mother-space norm is exactly $\norm g_G$.

The character identity
$U_aU_b=\sum_{k=|a-b|,\ k\equiv a+b\ (2)}^{a+b}U_k$
has $\min(a,b)+1$ terms. The weights of its terms are bounded by the
product weight. It follows that
\begin{equation}\label{eq:fusion}
 \norm{v_1v_2}_U\le
 \min\{\norm{v_1}_{U,1}\norm{v_2}_U,\,
       \norm{v_1}_U\norm{v_2}_{U,1}\}.
\end{equation}
An even Laurent multiplier with frequencies in $2m\Z$ has $U$-operator
norm at most its Wiener norm: use the two-character expression for a
paired cosine. In particular multiplication by $\abs p^2$ costs at most
$P^2$, where $P=\sum_jj\rho^{j-1}|c_j|$.

The absolute column sums in \eqref{eq:kd0}--\eqref{eq:kds} give
\begin{equation}\label{eq:kappa}
 \norm{K_D}_{G\to G}\le\kappa_m
 :=\max\left\{\frac1{2(m+1)(m+2)},\,\frac1{(m+2)(m+4)}\right\}.
\end{equation}
Thus $\kappa_1=1/12$ and $\kappa_6=1/80$. From
\eqref{eq:normal} and the shape weights,
\begin{align}
 \norm{q(g)}_U&\le\frac{\rho^{-m}}{2(m+1)}\norm g_G,&
 \norm{q(g)}_{U,1}&\le\frac{\rho^{-m}}{2m}\norm g_G,\label{eq:qbound}\\
 \sum_j\rho^{j-1}|\delta c_j|&\le
       \frac{\norm{\delta c}_C}{\sigma(1+2m)},&
 \sum_jj\rho^{j-1}|\delta c_j|&\le
       \frac{\norm{\delta c}_C}{\sigma}.\label{eq:thetacorrect}
\end{align}
These estimates and finite character fusion prove that \eqref{eq:map}
is a continuous polynomial $X\to\mathcal Y$.

No exponential radial weight is imposed on $s$. This causes no
regularity assumption on the solution: $|S_{n,s}|\le1$ on the disc, so
source truncations converge uniformly. The zero-Dirichlet
$W^{2,a}$ estimates for any finite $a>2$ give $C^1$ convergence of their
solutions and convergence of the normal traces. Consequently
$V=K_Dg$ has the stated boundary values and solves $\Delta V=g$ weakly.
At an exact zero the elliptic equation gives interior analyticity.
Analyticity across the boundary is addressed below.

\section{A contraction criterion}\label{sec:contraction}
Let $A:\mathcal Y\to X$ be a bounded injective operator. Let $x^0$ be
the centre used to construct $A$ and let $\bar x$ be the centre of the
contraction ball. Write
\[
 Y\ge\norm{A\F(\bar x)},\quad
 Z_0\ge\norm{I-AD\F(x^0)},\quad
 Z_1\ge\norm{A(D\F(\bar x)-D\F(x^0))}.
\]
Suppose that expanding the polynomial at $\bar x$ gives multilinear
remainder bounds $d_q$, $2\le q\le 2m+2$. Set $C_q=\norm A\,d_q$ and
\[
 R(r)=\sum_{q=2}^{2m+2}C_qr^q,\qquad
 Z_2(r)=\sum_{q=2}^{2m+2}qC_qr^{q-1}.
\]
Thus $Z_2$ denotes the nonlinear derivative contribution at radius $r$.
The inequalities
\begin{equation}\label{eq:radii}
 Y+(Z_0+Z_1-1)r+R(r)<0,\qquad
 Z_0+Z_1+Z_2(r)<1
\end{equation}
imply existence of a unique zero in $\overline B_X(\bar x,r)$.
Indeed $x\mapsto x-A\F(x)$ maps this complete closed ball into itself and
is a strict contraction. Its fixed point satisfies $A\F(x)=0$;
injectivity of $A$ gives $\F(x)=0$. The factor $q$ in $Z_2$ is obtained
by replacing one of the $q$ multilinear factors by a direction vector.

\section{The certificate in dimension four}\label{sec:r4}
\subsection{Finite centre and approximate inverse}
The centre $\bar x=x^0$ is given by
\path{finite_center_r4_M31_S24.json}. Its hexadecimal coefficients are
exact dyadic rationals. The source coordinates are $n=1,3,\ldots,31$,
$s=0,\ldots,24$, and the shape coordinates are $j=1,3,\ldots,31$.
There are $400+16=416$ finite unknowns. For orientation,
$b=c_1=7.08373029362\ldots$,
$c_3=0.153114245\ldots$ and $c_5=0.171014117\ldots$.

Let $\Pf$ and $\Pt$ project onto these finite coordinates and their
complements, in the appropriate source and boundary spaces.
The finite Jacobian $J_f=\Pf D\F(\bar x)|_{X_f}$ is assembled from four
disjoint interval column files. The verifier computes an interval
inverse, takes its exact dyadic midpoint $A_f$, and recomputes its
defect. Every midpoint entry is compared with the saved inverse.
The resulting bounds are
\begin{equation}\label{eq:r4finite}
 \norm{A_f}<1665.434,\qquad
 E_f:=\norm{I-A_fJ_f}<2.05\cdot10^{-34}<21\cdot10^{-35}.
\end{equation}
The $X$-norm of the last column of $A_f$, without the input row weight,
is $<54.778$. These are direct interval comparisons; the additional
\path{check_inequalities.py} also checks the sharper defect bound in
\eqref{eq:r4finite}. Thus $A_f$ and $J_f$ are invertible.

On high shape indices define the ball operator
\begin{equation}\label{eq:r4tail}
 (Tc)_j=-b^3(j+2)(c_j-c_{j+2}),\qquad
 (T^{-1}r)_j=-\sum_{\substack{k\ge j\\k\ {\rm odd}}}
                    \frac{r_k}{b^3(k+2)},\quad j\ge33.
\end{equation}
The ball shape column has coefficients $-b^3(j+2)$ at character row $j$
and $b^3j$ at row $j-2$, by differentiating \eqref{eq:map} at
$\psi=bw$, $h=b$, $q=b^2$. Weighted summation of
\eqref{eq:r4tail} gives
\[
 \norm{T^{-1}}\le L_T=(1-\rho^{-2})^{-1}.
\]
For $y=(f,v)\in\mathcal Y$, set $g_t=f_t$ and
\begin{equation}\label{eq:r4A}
 c_t=T^{-1}\Pt\bigl(v-2q_0q(f)\bigr),\qquad q_0=q(g^0).
\end{equation}
Subtract $b^3\,33\,c_{33,t}$ from row $31$ of the finite boundary output
$\Pf y$, and apply $A_f$ to that corrected finite output. Together with
$(g_t,c_t)$ this defines $Ay$. The finite row correction is the ball
coupling from shape index $33$; it is part of the operator.

The operator is injective. If $Ay=0$, its high source component gives
$f_t=0$, its high shape component gives
$\Pt(v-2q_0q(f))=0$, and the invertible finite map gives
$f_f=v_f=0$, since the row correction is then zero. Thus $f=0$, $v_t=0$
and $y=0$.
Put $N_q=\norm{q_0}_{U,1}<86.059$ and
$E=54.778(33/35)\rho^{-32}$. Equations
\eqref{eq:qbound}, \eqref{eq:r4tail} and the finite column bound give
\begin{equation}\label{eq:r4Anorm}
 \norm A\le1665.434+1+(1+E)L_T\frac{N_q}{2\rho}<2444.
\end{equation}
The complete residual is expanded without discarding any output
coefficient and preconditioned with \eqref{eq:r4A}. Its finite and tail
norms are respectively $<1.198868\cdot10^{-6}$ and
$<1.838234\cdot10^{-6}$, giving
\begin{equation}\label{eq:r4Y}
 Y<3.038\cdot10^{-6}.
\end{equation}

\subsection{All derivative columns}
In a weighted $\ell^1$ space an operator norm is the supremum of its
normalized column norms. The column partition is as follows:
\begin{center}
\begin{tabular}{@{}p{0.18\textwidth}p{0.59\textwidth}r@{}}
\toprule
Class & Indices & Upper bound\\
\midrule
Finite & All $416$ finite coordinates & $0.578+E_f$\\
Near source & $n\le31$, $25\le s\le55$; and
 $33\le n\le61$, $0\le s\le55$; $1336$ columns & $0.679$\\
Far source & $n\ge63$ or $s\ge56$ & $0.886$\\
Near shape & $j=33,35,\ldots,299$; $134$ columns & $0.720$\\
Far shape & Every odd $j\ge301$ & $0.925823$\\
\bottomrule
\end{tabular}
\end{center}
The finite-tail, near-source and near-shape files enclose the whole
output of \eqref{eq:r4A}, including the finite row correction. Their
spans are checked for exact coverage, and all their columns were
regenerated in the complete replay.

Here are the analytic arguments for the unenumerated classes.
Let $P=\sum_jj\rho^{j-1}|c_j|<9.09$ and
$d=|p_0|^2K_De_{n,s}$ for a source column. Its source derivative is
$e_{n,s}+d$, and its boundary derivative is exactly $2q_0q(e_{n,s})$.
The high boundary residual in \eqref{eq:r4A} is therefore
$-2\Pt(q_0q(d))$. Multiplication by $|p_0|^2$ shifts angular degree by
at most $30$ and lowers the radial index by at most $30$.
$K_D$ lowers it by at most one more. Hence $n\ge63$ cannot reach the
finite angular rows, and $s\ge56$ cannot reach the finite radial rows.
For $n\ge63$ the trace quotient costs at most
$[2\rho(n-29)]^{-1}$. Its whole-column bound is
\begin{equation}\label{eq:r4farsource}
 \frac{P^2}{65\cdot67}
 \left(1+(1+E)L_T\frac{N_q}{34\rho}\right)<0.885833.
\end{equation}
For $s\ge56$, $n\le61$, the harmonic row of $d$ vanishes, and the bound
is $P^2/(113\cdot115)<0.006359$. These cases cover the complement of the
near-source list.

For a high shape column, the pure boundary output, aligned by offsets
from $j$, has the form
\begin{equation}\label{eq:affine}
 B^{\rm shape}_{j+a}=jA_a+B_a,\qquad
 a=-92,-90,\ldots,60,\qquad \sum_aA_a=0,\quad j\ge125.
\end{equation}
To verify the identity analytically, expand
$-2h_0\,\chi_j|p_0|^2-h_0^2
(jw^{j-1}\bar p_0+j\bar w^{j-1}p_0)$ in characters.
The centre has character degree at most $30$ and derivative frequency
at most $30$. Once the high frequency is separated from these supports,
character fusion has fixed offsets and coefficients independent of $j$,
apart from the displayed derivative factor $j$. The latter is a
difference of adjacent characters, using
$2\cos(k\theta)=U_k-U_{k-2}$. The two fused sums have equal numbers of
terms, so their coefficient sums cancel. This proves the last identity
in \eqref{eq:affine}. The two exact columns $j=125,127$ determine
$A_a,B_a$; the program also checks $j=129,301$ exactly.

Put $D_a=B_a-(a+2)A_a$. Formula \eqref{eq:r4tail} gives
\[
 c_{j+v}=-b^{-3}\sum_{a\ge v}
       \left(A_a+\frac{D_a}{j+a+2}\right).
\]
For offsets $v<-92$ the $A_a$ sum is zero. Let
\begin{align*}
 L_v&=-b^{-3}\sum_{a\ge v}A_a-\mathbf1_{\{v=0\}},&
 L_0^\sharp&=\sum_v\rho^v|L_v|,\\
 L_1^\sharp&=\sum_v|v|\rho^v|L_v|,&
 D^\sharp&=b^{-3}\sum_v\rho^v\sum_{a\ge v}|D_a|,\\
 D^-&=\frac{\rho^{-94}}{b^3(1-\rho^{-2})}\sum_a|D_a|.
\end{align*}
The sums over $v$ here run through the offsets in \eqref{eq:affine}.
At $J=301$, the normalized boundary defect is bounded, for every
$j\ge J$, by
\begin{equation}\label{eq:r4boundaryfar}
 L_0^\sharp+\frac{L_1^\sharp}{J+2}
 +\left(1+\frac{60}{J+2}\right)\frac{D^\sharp}{J-90}
 +\frac{D^-}{J-90}<0.706765.
\end{equation}
This follows by bounding $(j+v+2)/(j+2)$ in the finite sums and extending
the low-offset leakage to its geometric series.

The interior shape output uses the zero-Dirichlet factor
$V_0=(1-r^2)Q_0$. The finite rational identity is checked coefficient by
coefficient, independently of its interval evaluation.
For angular degree $n>0$ the $r^2$ recurrence has coefficients
\[
 a_s=\frac{(s+1)(n+s+1)}{(n+2s+1)(n+2s+2)},\qquad
 c_s=\frac{s(n+s)}{(n+2s)(n+2s+1)},\qquad b_s=1-a_s-c_s.
\]
Thus $a_s,c_s\ge0$ and
$a_s+c_s\le(2s+1)/(n+1)$. The absolute column sum of $1-r^2$ is
$2(a_s+c_s)$.
Set $W=\bar p_0Q_0$, expanded in the complex mother basis.
Its least angular frequency is $-61$. Multiplication by $w^{j-1}$
increases the radial index by at most $\max(0,-k)$ for an input
frequency $k$, while its coefficients remain nonnegative with sum one.
The two conjugate shape terms consequently have normalized source norm
\begin{equation}\label{eq:r4interiorfar}
 I_J=\frac4{b^3}\sum_{k,s}
       \frac{|W_{k,s}|\rho^k[2(s+\max(0,-k))+1]}{J+k}
       <0.076108.
\end{equation}
The additional tail shape norm from its trace is at most
$L_TN_q I_J/[\rho(J-61)]<0.142951$.
The only finite output is the row correction in \eqref{eq:r4A}.
The boundary and trace contributions to it are bounded by
\[
 \frac{54.778\,33\sum_a|D_a|}
 {b^3(J-90)(J+2)\rho^{J-1}}
 +\frac{54.778\,33\,\rho^{-(J-93)}}{J-90}
                  \frac{N_q I_J}{\rho(J-61)}
 <5.23\cdot10^{-10}.
\]
Together with \eqref{eq:r4boundaryfar} and \eqref{eq:r4interiorfar}
this gives the far-shape bound $0.925823$. Every denominator increases
with $j\ge J$ and every leakage weight decreases. The estimate therefore
applies to the entire tail.

The five classes now give
\begin{equation}\label{eq:r4Z}
 Z_0=\norm{I-AD\F(\bar x)}<463/500=0.926,\qquad Z_1=0.
\end{equation}

\subsection{Nonlinear bounds and radius}
At the centre put
$P=\norm{p_0}_{\mathrm W}$, $P_2=\norm{|p_0|^2}_{\mathrm W}$,
$H_0=\norm{h_0}_U$, $H_1=\norm{h_0}_{U,1}$,
$H_{20}=\norm{h_0^2}_U$, $V_0^\sharp=\norm{K_Dg^0}_G$.
For a unit $X$ perturbation the linear bounds are
\[
 a_p=a_{h,1}=b^{-3},\quad a_{h,0}=(3b^3)^{-1},\quad
 a_{q,0}=(4\rho)^{-1},\quad a_{q,1}=(2\rho)^{-1},\quad K=1/12.
\]
Let $L_h=2\min(H_1a_{h,0},H_0a_{h,1})$ and
$Q_h=a_{h,0}a_{h,1}$. Expansion of \eqref{eq:map}, using
\eqref{eq:fusion}, gives
\begin{align*}
 d_{2,I}&=2Pa_pK+V_0^\sharp a_p^2,&d_{3,I}&=Ka_p^2,\\
 d_{2,B}&=a_{q,0}a_{q,1}+H_{20}a_p^2+2Pa_pL_h+Q_hP_2,\\
 d_{3,B}&=L_ha_p^2+2Pa_pQ_h,&d_{4,B}&=Q_ha_p^2.
\end{align*}
The interval computation of these expressions and \eqref{eq:r4Anorm}
gives
\[
 C_2<268,\qquad C_3<23/10000,\qquad C_4<6/10^8.
\]
At the exact rational radius $r=1/12500$, the rational upper bounds in
\eqref{eq:r4Y} and \eqref{eq:r4Z} give
\begin{align}
 Y+(Z_0-1)r+268r^2+\frac{23}{10000}r^3+\frac6{10^8}r^4
     &<-1.1667\cdot10^{-6},\label{eq:r4poly}\\
 Z_0-1+536r+\frac{69}{10000}r^2+\frac{24}{10^8}r^3
     &<-0.0311.\label{eq:r4lip}
\end{align}
These are strict rational comparisons. Section~\ref{sec:contraction}
gives an exact zero $x^*$ in this ball.

\subsection{Geometric and sign margins}
For a perturbation in this ball,
\[
 |\delta c_1|\le\frac r{3b^3},\qquad
 \sum_{j\ge1}j|\delta c_j|R^{j-1}\le\frac r{b^3}
 \quad(R\le\rho),\qquad
 \sum_{j\ge3}j(j-1)|\delta c_j|
 \le\frac{r\rho}{b^3e\log\rho}.
\]
The last estimate uses $\sup_{t\ge0}t\rho^{-t}=1/(e\log\rho)$.
The interval sums at the centre, with these perturbation bounds, give
\begin{align}
 \re\psi'(w)&>5.39635 &&(|w|\le21/20),\label{eq:r4geom1}\\
 |\psi'(w)|&>5.66945,\quad
 \re\left(1+\frac{w\psi''(w)}{\psi'(w)}\right)>0.10935
                                      &&(|w|\le1),\label{eq:r4geom2}\\
 \re\frac{w\psi'(w)}{\psi(w)}&>0.84029 &&(|w|=1),&
 |c_3^*|&>0.1531.\label{eq:r4geom3}
\end{align}
For example the convexity lower bound is
$1-(\sum_{j>1}j(j-1)|c_j|+r\rho/(b^3e\log\rho))/
(b-\sum_{j>1}j|c_j|-r/b^3)$.
The positive real derivative gives univalence by integrating along
line segments. The curvature and radial inequalities give strict
convexity and strict star-shapedness.

The character bound $|U_{n-1}|\le n$ gives $q^*>24.345$.
The same estimate for $h$ gives
$h^*\ge b-\sum_{j>1}j|c_j|-r/b^3>5.66945$; thus its positivity is also
a consequence of the derivative margin. The squared boundary equation
therefore yields the positive branch $q^*=h^*|p^*|$.
The interval values at $w=i/4,3i/4$, including the uniform field error
$r/12$, satisfy
\begin{equation}\label{eq:r4sign}
 V^*(i/4)>10,\qquad V^*(3i/4)<-8.
\end{equation}
The centre values are enclosed near $13.53518938515$ and
$-8.60629803547$, respectively. The two images lie on the positive
imaginary axis, since $\re\psi'>0$. Thus division by $y$ preserves
their opposite signs.

\section{The certificate in dimension fourteen}\label{sec:r14}
\subsection{Two centres and the finite block}
All norms in this section use $\rho=51/50$. The inverse centre $x^0$
is specified by \path{center_r14_M114_S36_arbrefined.json}, and the
contraction centre $\bar x$ by
\path{center_r14_M186_S48_arbrefined2.json}. They have the same exact
dyadic field scale $B$. At $x^0$ the finite source indices are
$n=6,18,\ldots,114$, $s=0,\ldots,36$, and the finite shape indices are
$j=1,13,\ldots,109$. There are $370+10=380$ finite unknowns. At $\bar x$
the corresponding cutoffs are $186$ and $48$.

The saved binary64 matrix $A_f$ is interpreted entrywise as an exact
dyadic matrix. Eight disjoint files enclose all finite Jacobian columns
at $128$-bit precision. Direct interval multiplication gives
\begin{equation}\label{eq:r14finite}
 \norm{A_f}<82859,\qquad
 \norm{I-A_fJ_f}<2.161\cdot10^{-9}<3\cdot10^{-9}.
\end{equation}
The program also stores and recomputes the norm of each output column
of $A_f$, including its row weight. This profile is used for the source
cross bound. Both $A_f$ and $J_f$ are invertible.

\subsection{The tail Schur operator}
At $x^0$ write
\[
 Qf=\frac{2q_0q(f)}b,\qquad
 H=D_c\F_1(x^0),\qquad \mathcal T=D_c\F_2(x^0),\qquad
 C=\Pt(\mathcal T-QH)|_{C_t},\quad b=c_1^0.
\]
The projections onto the high shape and boundary coordinates are
identified by boundary index $n=j+5$. Put
$h_b=b^6/B$, $\alpha=h_b^2$ and $\tau=(1-\rho^{-12})^{-1}$.
The ball shape operator has columns
\[
 \mathcal T_0e_j=-\alpha\bigl[(j+12)\chi_{j+5}-j\chi_{j-7}\bigr].
\]
Its step-$12$ tail inverse satisfies
$\norm{\mathcal T_0^{-1}}\le\tau/\alpha$, by the same weighted summation as
\eqref{eq:r4tail}. Define $\mathcal K=I-\mathcal T_0^{-1}C$.
The $20$ columns $j=121,133,\ldots,349$ are enclosed separately and have
maximum $<0.059850$. We next bound every $j\ge J=361$ analytically.

Put $P=\sum_jj\rho^{j-1}|c_j^0|$,
$C_\psi=\sum_j\rho^{j-1}|c_j^0|$,
$\delta P=\sum_{j>1}j\rho^{j-1}|c_j^0|$,
$\delta h_1=\norm{h_0-h_b}_{U,1}$ and
$\delta h_{20}=\norm{h_0^2-\alpha}_U$. Set
\[
 D_h=\frac{6C_\psi^5}{B},\quad
 \delta D_h=\frac{6(C_\psi^5-b^5)}B,\quad
 D_{h,b}=\frac{6h_b}b,\quad
 \delta P_2=2b\,\delta P+(\delta P)^2.
\]
Here $C_\psi\ge b$, and the nonnegative power expansion bounds the
coefficient difference from the ball by $\delta D_h$.
Expansion of $-2h_0\delta h\,|p_0|^2-h_0^2\delta|p|^2$, followed by
\eqref{eq:fusion} and division by the shape weight, bounds $\mathcal T-\mathcal T_0$ by
\[
 B_J=\frac2{b(J+12)}
 \bigl(\delta h_1D_hP^2+h_b\delta D_hP^2+
       h_bD_{h,b}\delta P_2\bigr)
       +\frac2b(P\delta h_{20}+\alpha\delta P).
\]
After $\mathcal T_0^{-1}$ this contributes $\tau B_J/\alpha<0.275974$.

For the other Schur term factor $K_Dg^0=(1-r^2)Q_0$ and set
$W=\bar p_0Q_0$. The factor identity is checked over $\mathbb Q$;
its interval expansion satisfies the same identity.
For a mother-space term at angular frequency $k$, multiplying by
$w^{j-1}$ gives frequency $j-1+k>0$ for $j\ge J$.
The nonnegative multiplication coefficients sum to one. In the
harmonic output row of $(1-r^2)\Phi_{n,s}$ only $s=0,1$ can contribute.
Their coefficients have magnitude at most $1/(n+2)$; taking the normal
trace costs a further $1/[2(n+1)]$. Consequently, with
$N_q=\norm{q_0}_{U,1}$, the $QH$ contribution to $\mathcal K$ is at most
\begin{equation}\label{eq:r14fartrace}
 \frac{4\tau N_q}{\alpha b\rho^6}
 \sum_{k,s}\frac{|W_{k,s}|\rho^k}{(J+k)(J+k+1)}
 <0.197588.
\end{equation}
This retains the harmonic row, rather than replacing it by the entire
$K_D$ column sum. Both denominator factors increase with $j\ge J$.
The total far bound is $<0.473562$. Combining near and far columns gives
\begin{equation}\label{eq:r14K}
 \norm{\mathcal K}<0.474,\qquad
 \norm{C^{-1}}\le\frac{\tau}{\alpha(1-\norm{\mathcal K})}<9.
\end{equation}
This covers the whole high shape space.

We use the finite polynomial
\begin{equation}\label{eq:CN}
 C_{13}^{-1}=\left(\sum_{\ell=0}^{12}\mathcal K^\ell\right)\mathcal T_0^{-1}.
\end{equation}
It has norm $<9$ and is invertible: its polynomial factor satisfies
$(\sum_{\ell=0}^{12}\mathcal K^\ell)(I-\mathcal K)=I-\mathcal K^{13}$, and both $I-\mathcal K$ and
$I-\mathcal K^{13}$ are invertible since $\norm{\mathcal K}<1$.

\subsection{The full approximate inverse and cross bounds}
For $y=(f,v)\in\mathcal Y$, define
\begin{align}
 \widetilde v&=v-Qf,&
 c_t&=C_{13}^{-1}\Pt\widetilde v,&
 g_t&=\Pt f-\Pt Hc_t,\label{eq:r14A1}\\
 x_f&=A_f\Pf\bigl(y-D\F(x^0)(g_t,c_t)\bigr),&
 Ay&=(x_f,g_t,c_t).\label{eq:r14A2}
\end{align}
The output shear, the invertible high shape map, the source shear and
the finite invertible map are triangular factors. More directly,
$Ay=0$ implies $c_t=g_t=0$, hence $\Pt f=0$ and
$\Pt(v-Qf)=0$, and then $x_f=0$ implies $\Pf f=\Pf v=0$.
It follows that $y=0$. The same equations reconstruct $y$ from any
$x_f,g_t,c_t$, so $A$ is also a bounded bijection.

The cross bounds are
\begin{equation}\label{eq:r14cross}
 X_s:=\norm{A_f\Pf D_g\F|_{G_t}}<0.171168<1/5,\qquad
 X_c:=\norm{A_f\Pf D_c\F|_{C_t}}<9.183045<10.
\end{equation}
We describe their complete support coverage. Let $p_{\max}=108$,
$h_{\max}=6\cdot109=654$ and $h^2_{\max}=2h_{\max}-6=1302$.
For the source bound, a pair of derivative monomials
$w^a\bar w^b$ is bounded by weight $\rho^{a+b}$ and has angular
displacement at most $a+b$, while its radial lowering is at most
$a+b$. Multiplication positivity and the exact $K_D$ column sum
bound each finite output by the largest inverse-profile entry among
the allowed output rows. The harmonic boundary output is expanded by
character fusion. Enumerating $n\le330$, $s\le253$, outside the finite
block, gives $6742$ possible source columns. For every remaining column
the finite source output is zero by support, and the finite boundary
output is zero either because $s>0$ or because the character supports
are separated. The certified maximum in \eqref{eq:r14cross} is reached
at $(n,s)=(6,37)$.

For the shape bound the finite output is exactly zero when
\[
 j+5-h_{\max}-p_{\max}>114,\quad
 j-1-p_{\max}-h^2_{\max}>114,\quad
 j-1-114-p_{\max}>114.
\]
All three hold for $j\ge1537$. The $118$ columns
$j=121,133,\ldots,1525$ are expanded and multiplied by $A_f$ with
interval arithmetic. Their index lists and centre and inverse hashes
are checked for complete coverage. Thus neither cross bound is a
sampled tail estimate.

The source smoothing and high shape derivative bounds are
\begin{equation}\label{eq:r14smallblocks}
 \norm{\Pt|p_0|^2K_D|_{G_t}}<0.141426<0.142,\qquad
 \norm{H|_{C_t}}<250.
\end{equation}
For the source estimate use $P^2$ times the largest high-source $K_D$
column sum: compare the lowest high radial index $s=37$ and the lowest
high angular index $n=126$, including both harmonic and nonharmonic
columns. For the shape estimate the $20$ near columns are enclosed
separately. The far bound comes from the factor $1-r^2$, as in
\eqref{eq:r4interiorfar}, and equals
\[
 4\sum_{k,s}\frac{|W_{k,s}|\rho^k
                 [2(s+\max(0,-k))+1]}{361+k}<249.711<250.
\]

For the feedback block, retain only those terms in \eqref{eq:kd0} and
\eqref{eq:kds} that can reach radial row zero after multiplication.
If a multiplier pair has total degree $d=a+b$, only $K_D$ terms at
radial index at most $d$ contribute, and their resulting angular degree
is at least $\max(6,n-d)$. The normal quotient therefore has the
additional factor $1/[2(\max(6,n-d)+1)]$. Summing these nonnegative
majorants proves
\begin{equation}\label{eq:r14QD}
 \norm{Q|p_0|^2K_D|_{G_t}}<0.000112959<114/10^6.
\end{equation}
The program checks $5952$ cases with $n\le342$, $s\le217$, outside the
finite block. For $s\ge218$ no harmonic output is possible; for
$n>342$ both the $K_D$ coefficients and the angular trace factors
decrease. This proves the bound on the unenumerated complement.

For the norm of $A$ on the entire output space, the coarse estimates
$\norm H<1076$ and $\norm Q<0.133<1$ follow from
$2P\norm{K_Dg^0}_G$ and \eqref{eq:qbound}.
Equations \eqref{eq:r14A1}--\eqref{eq:r14A2} yield
\begin{equation}\label{eq:r14Anorm}
 \norm A\le82859+(1+X_s)
       +\bigl((1+X_s)\norm H+X_c+1\bigr)\,9
       <94580<10^5,
\end{equation}
where $X_s=1/5$, $X_c=10$ may be used as rational ceilings.

\subsection{Every column of the derivative defect}
For a normalized high source input, the transformed boundary output in
\eqref{eq:r14A1} is $-Q|p_0|^2K_De$. Its shape correction has norm
at most $c=9(114/10^6)=513/500000$. The source error has norm at most
$g=0.142+250c=797/2000$. The finite correction is bounded by
$X_sg+X_cc$. Thus the whole column defect is bounded by
\begin{equation}\label{eq:r14highsource}
 (1+X_s)g+(1+X_c)c=244743/500000=0.489486<0.49.
\end{equation}
For a high shape input, \eqref{eq:CN} gives the exact shape remainder
$\mathcal K^{13}e$. Its source remainder is $H \mathcal K^{13}e$, and its finite remainder
is controlled by the two cross blocks. Consequently every such column
has defect at most
\begin{equation}\label{eq:r14highshape}
 \bigl((1+X_s)250+1+X_c\bigr)(0.474)^{13}<0.018962.
\end{equation}

For a finite normalized input, let $(f,v)=D\F(x^0)e$ and set
$c_0=\mathcal T_0^{-1}\Pt(v-Qf)$. The same algebra gives the coarse bound
\begin{equation}\label{eq:finitecoarse}
 E_f+\frac65\norm{\Pt f}
       +\frac{311}{1-0.474}\norm{c_0},\qquad E_f=3\cdot10^{-9}.
\end{equation}
Eight interval files cover all $380$ finite input columns.
For indices $0,\ldots,378$, \eqref{eq:finitecoarse} is $<0.114825$.
The final shape column, index $379$, has a larger coarse estimate and is
treated directly.

For that column let $\hat c$ be the eight-term dyadic shape vector in
\path{r14_M114_lastshape_tail_seed.json} and put
$R=\Pt(v-Qf)-C\hat c$.
Full coefficient expansion gives $\norm R<3.573\cdot10^{-15}$.
The difference between $C_{13}^{-1}\Pt(v-Qf)$ and $\hat c$ is at most
\[
 \epsilon_c=9\norm R+
       \frac{0.474^{13}}{1-0.474}\norm{c_0}.
\]
The second term compares the finite polynomial with the true Schur
inverse. Let $\hat g=\Pt f-\Pt H\hat c$. Expand its complete finite
correction by \eqref{eq:r14A2}. The sum of its finite, source-tail and
shape-tail defects, plus $311\epsilon_c$ divided by the input weight,
is $<0.160581$. Thus index $379$ is included with a rigorous remainder,
rather than accepting its finite seed as an exact solution.

Finite, high source and high shape inputs exhaust $X$. The preceding
estimates establish
\begin{equation}\label{eq:r14Z0}
 Z_0=\norm{I-AD\F(x^0)}<0.49.
\end{equation}

\subsection{Nonlinear bounds and centre transfer}
For a unit shape perturbation,
\[
 \theta_0=1/13,\qquad
 \theta_1=\sup_{j\equiv1\ (12)}\frac{j}{j+12}=1
\]
bound $\sum\rho^{j-1}|\delta c_j|$ and
$\sum j\rho^{j-1}|\delta c_j|$, respectively. Both estimates use the weighted norms in \eqref{eq:thetacorrect}.

At either centre let $H_{0,0}=\norm h_U$, $H_{1,0}=\norm h_{U,1}$.
For $1\le k\le6$ define
\begin{align*}
 H_{0,k}&=\frac{\binom6k}{B}C_\psi^{6-k}\theta_0^k,\\
 H_{1,k}&=\frac{\binom6k}{B}
 \left(\frac{6-k}{6}P C_\psi^{5-k}\theta_0^k
       +\frac{k}{6}C_\psi^{6-k}\theta_1\theta_0^{k-1}\right).
\end{align*}
The term with $6-k$ is zero when $k=6$.
These follow by expanding $\psi^6$ and differentiating a product of six
holomorphic factors; the character multiplier $n/6$ accounts for the
total degree divided by six. Put
\[
 L_\ell=\sum_{\substack{0\le k,k'\le6\\k+k'=\ell}}
       \min(H_{1,k}H_{0,k'},H_{0,k}H_{1,k'}),
 \quad a_{q,0}=\rho^{-6}/14,\quad a_{q,1}=\rho^{-6}/12.
\]
Take $L_\ell=0$ outside $0\le\ell\le12$. The boundary majorant at degree
$q$ is
\[
 d_{q,B}=\frac1{c_1}
   \left(P^2L_q+2P\theta_1L_{q-1}+\theta_1^2L_{q-2}
             +\mathbf1_{\{q=2\}}a_{q,0}a_{q,1}\right).
\]
For the interior part,
$d_{2,I}=2P\theta_1\kappa_6+\norm{K_Dg}_G\theta_1^2$,
$d_{3,I}=\kappa_6\theta_1^2$, and the other interior coefficients are
zero. Set $d_q=d_{q,I}+d_{q,B}$, $2\le q\le14$.
These are multilinear bounds in the sum input norm.

At $x^0$, $\sum_{q=2}^{14}qd_q<37.149$. The complete weighted
distance between the centres is
$\delta=\norm{\bar x-x^0}<6.440\cdot10^{-18}<1$.
Therefore
\begin{equation}\label{eq:r14transfer}
\begin{aligned}
 Z_1&\le\norm A\sum_{q=2}^{14}qd_q\delta^{q-1}\\
    &<10^5\cdot37.149\cdot6.440\cdot10^{-18}<2.4\cdot10^{-11},\\
 Z_0+Z_1&<0.490000000024<1/2.
\end{aligned}
\end{equation}
This uses $\delta^{q-1}\le\delta$ for $q\ge2$.

At $\bar x$ the complete coefficient residual is
$<1.063\cdot10^{-30}$. The evaluator uses the denominator $\bar c_1$
in the boundary output, whereas \eqref{eq:map} fixes $b=c_1^0$.
Multiply its residual and nonlinear coefficients by
$f=\max(1,\bar c_1/b)<1+1.7\cdot10^{-23}$ to bound the fixed map.
Both nonlinear certificates use $\theta_1=1$, and their quadratic
coefficients satisfy $d_2<18.551$. Hence $Y<1.063\cdot10^{-25}$.

At $r=10^{-20}$, using $10^5$ and $1/2$ as ceilings for the inverse norm
and transferred defect, respectively, the interval comparisons give
\begin{align}
 Y-\frac r2+\sum_{q=2}^{14}10^5 f d_qr^q
        &<-4.99989\cdot10^{-21},\label{eq:r14poly}\\
 -\frac12+\sum_{q=2}^{14}q\,10^5 f d_qr^{q-1}
        &<-0.49999999999996.\label{eq:r14lip}
\end{align}
In particular $Z_2(r)<4\cdot10^{-14}$. The exact zero $x^*$ follows
from Section~\ref{sec:contraction}.

\subsection{Geometry and sign}
At $\bar x$ put $L=\bar c_1-\sum_{j>1}j|\bar c_j|$ and
$U=\sum_{j>1}j(j-1)|\bar c_j|$.
The radius bounds $|\delta\psi'|\le r$ on $|w|\le1.001$ and
$|\delta\psi''|\le50r$ on $|w|\le1$ follow from
\eqref{eq:thetacorrect} and $1/(e\log(51/50))<50$.
The centre sums certify a positive outer derivative lower bound and
$L-U-51r>25.4946>0$. Hence $\psi^*$ is univalent on
$|w|\le1.001$, and $\re(1+w\psi''/\psi')>0$ on the unit disc.
Its image is strictly convex and contains zero.
The interval comparison $|\bar c_{13}|-r/[25\rho^{12}]>0$
proves $c_{13}^*\ne0$.

For a character polynomial, its constant coefficient minus
$\sum_{n>6}(n/6)|v_n|$ is a lower bound on its boundary value.
The source perturbation costs at most $r\rho^{-6}/12$, and the shape
perturbation costs at most $\sum_{k=1}^6H_{1,k}r^k$.
The exact-boundary margins are
\begin{equation}\label{eq:r14positive}
 q^*>29.1875,\qquad h^*>0.98405.
\end{equation}
Thus the positive square root is obtained throughout the boundary.
At $w=t e^{i\pi/12}$ the centre values for $t=1/10,37/100$ are
enclosed below $-0.0795470$ and above $1.4919379$, respectively.
The field perturbation is at most $r/80$, so the exact values still
have opposite signs. Real coefficients $j=1\pmod{12}$ keep this ray
fixed, and its image has positive radius because $\re\psi'>0$.
Therefore $\pi(\psi(w))>0$ at both points.

\section{Reconstruction of the domains and fields}\label{sec:reconstruction}
\subsection{The planar field and exact boundary data}
At either exact zero, $V=K_Dg$ has zero Dirichlet trace and satisfies
$\Delta V+|p|^2V=0$. The conformal map is univalent on a neighbourhood
of the closed disc. Set $F=B V\circ\psi^{-1}$ on $D=\psi(\D)$.
The conformal normal derivative formula gives
\[
 \partial_\nu F(\psi(e^{i\theta}))
       =B\,\frac{Ng(\theta)}{|p(e^{i\theta})|}.
\]
The positive margins for $q,h$ give $Ng=h|p|\sin(m\theta)$.
Thus $\partial_\nu F=\im\psi$ for $m=1$ and
$\partial_\nu F=\im(\psi^6)=\pi$ for $m=6$.
These are equalities for the exact solution, including the wall
endpoints by continuity.

The field is analytic in $D$ by elliptic regularity. Its zero
Dirichlet data and the displayed analytic Neumann data on the analytic
boundary give a local analytic continuation by the
Cauchy--Kowalevski theorem. To identify it with $F$ on the interior
side, their difference has zero Cauchy data. Extend the difference by
zero across the boundary; integration by parts shows that it solves
Helmholtz distributionally there. Analytic elliptic regularity and
vanishing on the exterior open set imply that the difference is zero.
Compactness of the boundary gives an analytic continuation on a
neighbourhood of $\overline D$.

The real coefficient and sine symmetries commute with the equation
and the conformal change of variables. For $m=1$, $D$ is invariant
under both coordinate reflections and $F$ is odd in $y$.
For $m=6$, $\psi(e^{i\pi/6}w)=e^{i\pi/6}\psi(w)$ and
$F$ changes sign under the corresponding rotation. It is odd under
conjugation, and hence under the Weyl reflections. In particular $D$
is Weyl-invariant and $F$ is Weyl anti-invariant.

\subsection{The four-dimensional lift}
Define
\[
 \Omega_4=\{(x,X'):(x,|X'|)\in D\},\qquad
 u_4(x,X')=G(x,|X'|^2),\quad F(x,y)=yG(x,y^2).
\]
The division is analytic locally at every axis point. Away from the
axis it is ordinary division by $y$, so the local definitions agree.
Equation \eqref{eq:r4radial} proves the Helmholtz equation, and the
boundary product rule proves $u_4=0$, $\partial_\nu u_4=1$, extending
to the axis by analyticity.

Convexity follows from planar convexity and reflection symmetry:
for two points the norm of their transverse convex combination is at
most the corresponding convex combination of transverse norms.
The planar section contains every smaller nonnegative vertical
coordinate at the same horizontal coordinate. Strict convexity
follows as well. If the meridian data of two different boundary
points differ, their planar convex combination is interior. If they
coincide but the transverse directions differ, strict triangle
inequality decreases the transverse norm, again giving an interior
point. At transverse norm zero identical meridian data give a single
point, so this case has no exception.

For completeness the analytic boundary can be written as a radial
sphere. The positive meridian radius $R_D(\theta)$ is analytic,
even and $\pi$-periodic. Its exponentially convergent Fourier series
uses only $\cos(2k\theta)=T_k(\cos2\theta)$. On $S^3$ substitute
$\cos2\theta=\omega_1^2-|\omega'|^2$. The resulting positive function
$R(\omega)$ is analytic: the exponential coefficient bound makes the
Chebyshev series converge on a complex neighbourhood of the compact
range $[-1,1]$. The radial map $\omega\mapsto R(\omega)\omega$ is
injective, and its derivative has tangential projection
$R(\omega)$ times the identity. It is therefore an analytic
parametrization of the boundary by $S^3$.

The symmetries force the centre of a hypothetical ball to be zero.
Its meridian section would be a centred disc. A conformal map to that
disc with $\psi(0)=0$, $\psi'(0)>0$ is linear, contrary to
$c_3^*\ne0$. Finally \eqref{eq:r4sign}, on the positive meridian
half-plane, gives opposite signs of $u_4$. This proves
Theorem~\ref{thm:r4}.

\subsection{The invariant lift in dimension fourteen}
Set $\Omega_{14}=\Ad(G_2)D$. Every orbit meets $\tf$ in a Weyl orbit,
so $\Omega_{14}\cap\tf=D$ and the regular field is defined by
$u_{14}|_\tf=F/\pi$. Analytic anti-invariance makes $F$ vanish on each
wall. Analytic division by a linear wall factor can be repeated:
after division the remaining wall vanishing follows off the
intersections and then by analytic continuation. Thus $F/\pi$ is
analytic and Weyl-invariant on the whole Cartan section, including
the origin. The field scale and the normalization of $\pi$ have been
fixed throughout. Using the literal positive-root product instead
would multiply the constant flux by a nonzero scalar, removable by
scaling $u_{14}$.

We state the convexity result used here precisely. For a real
semisimple Cartan decomposition $\mathfrak g=\mathfrak t_c\oplus
\mathfrak p$, with maximal abelian section $\mathfrak a\subset
\mathfrak p$, Kostant's theorem says
$P_{\mathfrak a}(\Ad(K)x)=\operatorname{conv}(Wx)$ for
$x\in\mathfrak a$. Lewis~\cite{Lewis2000}, Theorems~2.2 and~3.6,
formulates its consequence: an $\Ad(K)$-invariant subset of
$\mathfrak p$ is convex exactly when its section in $\mathfrak a$
is convex. The projection theorem is due to
Kostant~\cite{Kostant1973}. Apply this to the complexification of
$\mathfrak g_2$ regarded as a real Lie algebra, with
$\mathfrak p=i\kf$, $\mathfrak a=i\tf$ and compact group $G_2$.
Multiplication by $i$ is an equivariant isometry after normalization
of the Killing form. The convex Weyl-invariant section therefore
gives convexity of $\Omega_{14}$.

One can also see the set conclusion from the projection statement.
The closed invariant lift of $\overline D$ equals the intersection
of all adjoint conjugates of
$\{X:P_\tf X\in\overline D\}$. The inclusion from left to right uses
$\operatorname{conv}(Wx)\subset\overline D$. For the reverse inclusion,
conjugate $X$ into $\tf$ and apply that particular projection.
The intersection is convex. Its interior is $\Ad(G_2)D$: the radial
description below identifies interior and boundary. Thus the open
domain is convex.

To justify the analytic boundary, let $R_D(\theta)$ be the positive
analytic radial function of $D$. Weyl symmetry makes it an analytic
function $f(\cos6\theta)$. At the endpoints of $[-1,1]$ this follows
from the even analytic Taylor expansion at a reflection wall and
the nonzero quadratic term of $\cos6\theta$ there. Polynomial
Chevalley restriction gives a real homogeneous invariant polynomial
$P_6$ on $\kf$ restricting to $\re z^6$; see, for example, the
restriction statement in~\cite{Tevelev1999}. Hence on an annulus
around the unit sphere
\[
 R(X)=f\left(\frac{P_6(X)}{|X|^6}\right)
\]
is analytic. Every unit direction is conjugate to a Cartan direction,
so its argument lies in $[-1,1]$ and $R$ is positive.
This yields an analytic radial boundary with nonzero normal,
parametrized by $S^{13}$.

The regular orbit field satisfies Helmholtz by
\eqref{eq:radialg2}. Its extension as an orbit function is locally
bounded since the analytic section quotient is locally bounded.
The nonregular set has codimension at least three: a nonzero wall
has one section parameter and an extra two-dimensional root plane
in its centralizer, so its orbit saturation has dimension
$1+(14-4)=11$; the origin has higher codimension. On compact subsets
away from the origin these are finitely many smooth wall strata.
An $\epsilon$-tube cutoff around a stratum of codimension $c\ge3$
has gradient and Hessian integrals
$O(\epsilon^{c-1})$ and $O(\epsilon^{c-2})$. Testing Helmholtz
against this cutoff, using boundedness of the field, makes all error
terms tend to zero. A separate ball cutoff treats the origin.
The equation consequently holds distributionally across the
nonregular set. Elliptic analytic regularity gives the analytic
extension. Apply the same argument to the planar continuation on
a slightly larger invariant neighbourhood to obtain boundary
regularity.

At regular boundary points $u_{14}=0$ and
$\partial_\nu u_{14}=1$ by \eqref{eq:fluxg2}. The regular boundary
points are dense, and the analytic field and boundary extend both
equalities to all boundary points. The domain is bounded, and its
Cartan section is strictly convex.
An adjoint-invariant ball would have an adjoint-fixed centre; the
simple Lie algebra has no nonzero fixed vector. Its section would
therefore be a centred disc, contradicting $c_{13}^*\ne0$.
The two points of Section~\ref{sec:r14} have positive $\pi$ and
opposite field signs. This proves Theorem~\ref{thm:r14}.

\section{The computer-assisted inequalities}\label{sec:computer}
The verification directories contain the two self-contained Python
programs, their transparent coefficient programs and exact data,
supplementary numerical comparisons, and SHA-256 lists. All
coefficient calculations use rational arithmetic or Arb interval
arithmetic through \texttt{python-flint 0.9.0}. The recorded Python
version is $3.10.19$ on Linux x86-64. Arb calculations use one thread;
the surrounding numerical libraries are also restricted to one
thread. Binary64 files specify exact dyadic inputs.

The four-dimensional single file authenticates its payload and
manifest, checks every finite and near-tail span, reconstructs and
compares the finite inverse, and recomputes the residual, far-tail,
nonlinear, rational-radius and geometric inequalities. Its complete
mode additionally regenerates every bounded column file and compares
the bytes. The fourteen-dimensional file authenticates its payload
and individual files, checks finite and cross-column coverage,
recomputes the inverse profile and far bounds, composes the
all-column defect, regenerates both nonlinear coefficient files
with $\theta_1=1$, transfers the centre, and checks the radius,
boundary signs and geometry. Its complete mode regenerates every
finite and near-column file.

The finite inverse and near columns do not themselves prove an
infinite-dimensional norm bound. The support arguments,
character identities, far-tail estimates, injectivity and
reconstruction are proved in the text. The programs enclose the
finite expressions in those estimates. The additional
\path{check_inequalities.py} compares their typed interval outputs
with the sharper numerical ceilings printed here. The finite matrix program enforces the ceiling $21\cdot10^{-35}$ used in the assembled estimate; the supplementary comparisons also
check the sharper finite defect in \eqref{eq:r4finite}.

The principal contraction bounds are:
\begin{center}
\begin{tabular}{@{}lcc@{}}
\toprule
Quantity & Dimension four & Dimension fourteen\\
\midrule
$r$ & $1/12500$ & $10^{-20}$\\
$Y$ & $<3.038\cdot10^{-6}$ & $<1.063\cdot10^{-25}$\\
$Z_0$ & $<0.926$ & $<0.49$\\
$Z_1$ & $0$ & $<2.4\cdot10^{-11}$\\
$Z_2(r)$ & $<0.042881$ & $<4\cdot10^{-14}$\\
$\norm A$ & $<2444$ & $<10^5$\\
\bottomrule
\end{tabular}
\end{center}

For location of the numerical comparisons, the finite inverse and
residual in dimension four are in
\path{verify_finite_inverse_r4.py}; its five column classes are in the
three bounded-column families and
\path{verify_far_source_bound_r4.py},
\path{verify_far_shape_bound_r4.py}. Its remaining inequalities are in
\path{verify_nonlinear_r4.py}, \path{verify_radii_r4.py},
\path{verify_geometry_r4.py} and \path{verify_signchange_r4.py}.
\path{verify_ball_linearization.py} also checks a Dirichlet-ball Bessel near-resonance and the quadratic and cubic moments of $U_4$; this is a consistency check on which no inequality of the proof depends.
In dimension fourteen the finite matrix is checked by
\path{verify_finite_block_rank2.py}; the tail and cross estimates are
checked by the Schur, source-tail, cross-source, cross-shape, $H$ and
$QD$ programs. The finite-input files and
\path{verify_last_finite_column_r14.py} complete that column class.
The nonlinear and centre-transfer calculations are in
\path{nonlinear_majorant_rank2.py} and
\path{verify_center_shift_r14.py}, and the final inequalities in
\path{verify_radii_geometry_r14.py}.

From the corresponding verification directory run, sequentially,
\begin{verbatim}
OMP_NUM_THREADS=1 OPENBLAS_NUM_THREADS=1 MKL_NUM_THREADS=1 \
  nice -n 19 python3 -B verify_berenstein_r4.py
OMP_NUM_THREADS=1 OPENBLAS_NUM_THREADS=1 MKL_NUM_THREADS=1 \
  nice -n 19 python3 -B verify_r14.py
\end{verbatim}
Each command ends with the line \texttt{PROVED} and exit code zero.
For complete regeneration append \texttt{--recompute} to the
four-dimensional command and \texttt{--recompute-all} to the
fourteen-dimensional command. Run the additional numerical comparisons
with \texttt{nice -n 19 python3 -B check\_inequalities.py}.
Its last line is \texttt{ALL\_PAPER\_INEQUALITIES\_CHECKED}.

Recorded single-file wall times are \rFourSeconds\ seconds and
\rFourteenSeconds\ seconds, respectively. Complete serial
regeneration of the supplied mathematical programs and numerical
data took \rFourFullSeconds\ seconds and
\rFourteenFullSeconds\ seconds.
Optimization with \texttt{-O} is rejected before mathematical
verification. In each dimension a changed embedded numerical file
was also rejected after its outer payload digest was recomputed,
because the inner file manifest did not match.

The single-file and compressed-payload SHA-256 values are
\noindent Dimension four, program:\par\nopagebreak
\noindent\texttt{\rFourHash}\par
\noindent Dimension four, payload:\par\nopagebreak
\noindent\texttt{\rFourPayload}\par
\noindent Dimension fourteen, program:\par\nopagebreak
\noindent\texttt{\rFourteenHash}\par
\noindent Dimension fourteen, payload:\par\nopagebreak
\noindent\texttt{\rFourteenPayload}\par
Each directory's \path{SHA256SUMS} also lists every transparent
verification file. A hash comparison authenticates the specified
reproducible inputs; the analytic proof and the interval calculations
establish the theorem.

\section{Remarks}
The two solutions lie outside the sign-definite case of Serrin's
theorem. In dimension fourteen the assertion of strict convexity is
made for the Cartan section; the ambient-domain assertion is convexity.
Dimensions six, eight and ten are not covered by these certificates.
The two coefficient balls specify the exact constructions.

\section*{Use of AI}
The proofs, the programs and the text of this paper were produced with
AI systems (GPT-6 Astra, GPT-6.1 Sol and Claude Opus~5.5) under the
author's direction. They were checked by independent reviews with
AI models and by independent replays of the interval-arithmetic
certificates. The author takes responsibility for the content.

\bibliographystyle{amsplain}
\bibliography{references}
\end{document}
